\documentclass[../../main.tex]{subfiles}% 注意这里的文件路径不能用 ./main.tex ,否则用latexmk编译子文件会报错
\graphicspath{{\subfix{./image/}}{\subfix{../../image/}}} % 指定图片引用路径,后续可以直接使用图片文件名
% 如果图片存放在上级目录,则使用 ../../image/ 作为备用路径

% 例如:
% \begin{figure}[H]
% \centering
% \includegraphics[width=0.8\textwidth]{图.png}
% \caption{}
% \label{figure:图}
% \end{figure}
% 注意:上述\label{}一定要放在\caption{}之后,否则引用图片序号会只会显示??.

\begin{document}

\section{圆规直尺作图}

圆规直尺作图指使用无刻度的直尺与圆规进行平面几何作图，该问题可通过\textbf{Galois理论}完全刻画可作图的几何图形。

\begin{lemma}
若长度$1,\alpha,\beta$已给出，则可用圆规和直尺作出$\frac{\beta}{\alpha},\sqrt{\alpha}$。
\end{lemma}

\begin{proof}\begin{enumerate}[(1)]
\item 作$\frac{\beta}{\alpha}$：过点$P$作直线$l_1,l_2$，在$l_1$上取点$Y,U$满足$PY=\alpha,\ YU=\beta$；在$l_2$上取点$E$满足$PE=1$。连接$YE$，过$U$作$YE$的平行线交$l_2$于$V$，则$EV=\frac{\beta}{\alpha}$。
\begin{figure}[H]
\centering
\includegraphics[scale=0.3]{fig6_6_1.png}
\caption{}
\label{figure:fig6_6_1}
\end{figure}

\item 作$\sqrt{\alpha}$：在直线上取点$P,E,F$满足$PE=1,\ EF=\alpha$。取$PF$中点$O$，以$O$为圆心、$\frac{1}{2}(1+\alpha)$为半径作圆弧，交过$E$的$PF$垂线于$G$，则$EG=\sqrt{\alpha}$。
\begin{figure}[H]
\centering
\includegraphics[scale=0.3]{fig6_6_2.png}
\caption{}
\label{figure:fig6_6_2}
\end{figure}
\end{enumerate}
同时可作出$\alpha\beta$。

\end{proof}

将平面几何作图问题代数化：在平面上建立复平面$\mathbb{C}$，平面点与复数一一对应。$\mathbb{Q}(\sqrt{-1})$中的元素均可尺规作出；若域$K$中元素可作出，则$K$的二次扩域$K(\gamma)$中元素也可作出。

尺规作图仅产生三类交点，对应域扩张性质：
\begin{enumerate}[(1)]
\item 直线与直线的交点：交点坐标属于已知点生成的域$K$，无域扩张；
\item 直线与圆的交点：交点坐标属于$K$的二次扩域$K(\gamma)$，满足$[K(\gamma):K]=2$；
\item 圆与圆的交点：两圆方程相减可化为直线与圆的交点问题，仍对应二次扩域。
\end{enumerate}
因此尺规作图仅能通过整数四则运算与有限次开平方得到新的长度。

\begin{theorem}
复数$\gamma\in\mathbb{C}$可用圆规直尺作出的充分必要条件是：存在二次扩域序列
\begin{align*}
\mathbb{Q}\subseteq K_1\subseteq K_2\subseteq\cdots\subseteq K_n,
\end{align*}
使得$\gamma\in K_n$。
\end{theorem}

\begin{proof}
\textbf{必要性}：由尺规作图的交点类型，每次作图仅产生二次扩域，故可作复数必属于有限次二次扩域链的末端域。

\textbf{充分性}：设$K$可作出，$E/K$为二次扩域，$\gamma\in E\setminus K$，其极小多项式为$\mathrm{Irr}(\gamma,K)=x^2+\alpha x+\beta,\ \alpha,\beta\in K$，则
\begin{align*}
\gamma=\frac{1}{2}\left(-\alpha\pm\sqrt{\alpha^2-4\beta}\right).
\end{align*}
由引理，$\sqrt{\alpha^2-4\beta}$可尺规作出，故$\gamma$可作出，$E$中所有元素均可作出。

\end{proof}

\begin{remark}
可进一步要求$K_n/\mathbb{Q}$为正规扩张。
\end{remark}

\begin{example}
尺规作图不能三等分任意角。
\end{example}

\begin{proof}
设需三等分角$3\theta$，由三角恒等式
\begin{align*}
4\cos^3\theta-3\cos\theta=\cos3\theta.
\end{align*}
令$x=\cos\theta,\ a=\cos3\theta$，则方程$4x^3-3x-a=0$的根一般无法通过有限次平方根表示，故任意角无法三等分。

若将无刻度直尺换为有刻度直尺，则可三等分任意角。
\begin{figure}[H]
\centering
\includegraphics[scale=0.3]{fig6_6_4.png}
\caption{}
\label{figure:fig6_6_4}
\end{figure}
设直尺端点为$O$，刻点为$P$，以$C$为圆心、$OP$为半径作圆，移动直尺使$P$在圆周上、直尺过$B$，则$\angle COP=\frac{1}{3}\angle ACB$。
由$OP=CP=CA$，得$\angle CBP=\angle CPB=2\angle COP$，结合三角形内角和可证该结论。

\end{proof}

\begin{example}
设$p$为素数，则正$p$边形可用圆规直尺作出当且仅当$p=2^{2^q}+1$（称$p$为\textbf{Fermat素数}）。
\end{example}

\begin{proof}
\textbf{必要性}：单位圆内接正$p$边形等价于作出$p$次本原单位根$\zeta$，分圆域满足$[\mathbb{Q}(\zeta):\mathbb{Q}]=p-1$。
由可作图条件，存在二次扩域链使得$\zeta\in K_n$，故$p-1=[\mathbb{Q}(\zeta):\mathbb{Q}]\mid[K_n:\mathbb{Q}]=2^n$，即$p-1=2^m$。
若$m$含奇因子$s>1$，则$p=(2^r)^s+1$可被$2^r+1$整除，与$p$为素数矛盾，故$m=2^q$，即$p=2^{2^q}+1$。

\textbf{充分性}：若$p=2^{2^q}+1$为素数，则分圆域的Galois群阶$|\mathrm{Gal}(\mathbb{Q}(\zeta)/\mathbb{Q})|=p-1=2^{2^q}$，为2的幂。
Galois群存在子群序列$G_0\supsetneq G_1\supsetneq\cdots\supsetneq G_n=\{\mathrm{id}\}$满足$[G_{i-1}:G_i]=2$，对应不变子域的二次扩域链$\mathbb{Q}=K_0\subseteq K_1\subseteq\cdots\subseteq K_n=\mathbb{Q}(\zeta)$，故$\zeta$可尺规作出。

\end{proof}

\begin{remark}
已知的Fermat素数为$2^{2^0}+1=3,\ 2^{2^1}+1=5,\ 2^{2^2}+1=17,\ 2^{2^3}+1=257,\ 2^{2^4}+1=65537$；$q\geqslant5$时$2^{2^q}+1$均为合数。
\end{remark}

\begin{example}
正十七边形（$17=2^{2^2}+1$）可通过尺规作图作出。
\end{example}
\begin{remark}
立方倍积问题、化圆为方问题均无法通过圆规直尺作出。
\end{remark}
\begin{proof}
取17次本原单位根$\xi=\cos\frac{2\pi}{17}+\sqrt{-1}\sin\frac{2\pi}{17}$，则$[\mathbb{Q}(\xi):\mathbb{Q}]=16$，$\mathrm{Gal}(\mathbb{Q}(\xi)/\mathbb{Q})=\langle\sigma\rangle$，其中$\sigma(\xi)=\xi^3$为16阶元。
Galois群子群序列：$\langle\sigma\rangle\supsetneq\langle\sigma^2\rangle\supsetneq\langle\sigma^4\rangle\supsetneq\langle\sigma^8\rangle\supsetneq\{\mathrm{id}\}$，对应不变子域链：
\begin{align*}
\mathbb{Q}\subseteq\mathbb{Q}(a_2)\subseteq\mathbb{Q}(a_2,a_4)\subseteq\mathbb{Q}(a_2,a_4,a_8)\subseteq\mathbb{Q}(\xi),
\end{align*}
其中
\begin{align*}
a_2&=\xi+\xi^2+\xi^4+\xi^8+\xi^9+\xi^{13}+\xi^{15}+\xi^{16},\\
a_4&=\xi+\xi^4+\xi^{13}+\xi^{16},\\
a_8&=\xi+\xi^{-1}=2\cos\frac{2\pi}{17}.
\end{align*}
各元素的极小多项式：
\begin{align*}
\mathrm{Irr}(a_2,\mathbb{Q})&=x^2+x-4,\quad a_2=\frac{1}{2}(-1+\sqrt{17}),\\
\mathrm{Irr}(a_4,\mathbb{Q}(a_2))&=x^2-a_2x-1,\quad a_4=\frac{1}{2}(a_2+\sqrt{8-a_2}),\\
\mathrm{Irr}(a_8,\mathbb{Q}(a_2,a_4))&=x^2-a_4x+b,\\
\mathrm{Irr}(\xi,\mathbb{Q}(a_2,a_4,a_8))&=x^2-a_8x+1,\quad \xi=\frac{1}{2}\left(a_8+\sqrt{a_8(a_8+4)}\right).
\end{align*}

\textbf{具体作图步骤}：
\begin{enumerate}[(1)]
\item 作单位圆$O$的两条垂直直径，得圆上点$P_0,B$；
\item 取$OJ=\frac{1}{4}OB$，构造$\angle OJE=\frac{1}{4}\angle OJP_0,\ \angle FJE=45^\circ$；
\item 以$FP_0$为直径作圆交$OB$于$K$，以$E$为圆心、$EK$为半径作圆交$OP_0$于$N_5,N_3$；
\item 过$N_5,N_3$作$OB$的平行线，交圆$O$于$P_3,P_5$，弧$\overset{\frown}{P_3P_5}$的平分点为$P_4$，则$P_3P_4$为正十七边形的边长。
\end{enumerate}
\begin{figure}[H]
\centering
\includegraphics[scale=0.3]{fig6_6_5.png}
\caption{正十七边形尺规作图示意}
\label{figure:fig6_6_5}
\end{figure}

\textbf{正确性验证}：只需证明
\begin{align}
OE=\frac{1}{2}\left(\cos\frac{3\times2\pi}{17}+\cos\frac{5\times2\pi}{17}\right),\label{eq:6_6_1}\\
ON_3=\cos\frac{3\times2\pi}{17}.\label{eq:6_6_2}
\end{align}
利用三角恒等式与域扩张的代数关系，可推得上述等式成立，故正十七边形可尺规作出。

\end{proof}





\end{document}